\section{Background and Related Work}
\label{sec:background}

\paragraph{Generative retrieval.} Generative retrieval replaces nearest-neighbor search with directly decoding an identifier for the target document. DSI~\cite{tay2022dsi} memorizes query-to-docid mappings in a Transformer's parameters; NCI~\cite{wang2022nci} and SEAL~\cite{bevilacqua2022seal} improve the identifier scheme via hierarchical clustering and document substrings; GENRE~\cite{decao2021genre} introduces constrained decoding over a trie of valid identifiers, adopted by nearly every later generative retriever including GENIUS. A parallel line replaces these with \emph{learned discrete semantic IDs}: TIGER~\cite{rajput2023tiger} applies RQ-VAE-style IDs to recommendation, GenRet~\cite{sun2023genret} learns the identifier scheme jointly with retrieval, and IRGen~\cite{zhang2023irgen} applies a similar idea to image retrieval; effectiveness degrades sharply with corpus size and task diversity~\cite{pradeep2023scale} (Section~\ref{sec:setup}; survey:~\cite{li2024survey}). GENIUS~\cite{kim2025genius} is the multimodal instance of this line. None of this prior work studies whether a semantic-ID codebook's balance regularization helps or hurts \emph{retrieval} -- the gap this paper closes.

\paragraph{RQ and the GENIUS pipeline.} An RQ turns a continuous embedding into discrete codes by quantizing in stages~\cite{lee2022rqvae}: each level's codebook finds the nearest of $K$ reference vectors to the residual left by the levels above it. GENIUS~\cite{kim2025genius} applies this to multimodal retrieval over M-BEIR~\cite{wei2024uniir}: a CLIP-style~\cite{radford2021clip} encoder (CLIP-SF) produces a joint image/text embedding; an 8-level RQ (plus a modality-indicator level) quantizes it into a semantic ID; a T5~\cite{raffel2020t5} decoder generates a target ID via constrained beam search over a trie of valid candidate IDs, following~\cite{decao2021genre} -- a prefix tree containing exactly the sequences that name a real candidate, so decoding can never emit a nonexistent ID. Each codebook is fit via an EMA update -- a running average of the embeddings assigned to it, rather than a gradient-trained parameter -- to whatever vectors it observes, with no term rewarding even code usage.

\paragraph{Codebook balance and collapse.} Vector-quantized models suffer \emph{codebook collapse}: a small code subset absorbs most assignments while the rest receive vanishing gradient~\cite{vandenoord2017vqvae}; the standard remedy adds an entropy-style auxiliary loss alongside reconstruction~\cite{zhang2023regularizedvq} -- \emph{balance regularization}. An alternative line avoids collapse by construction rather than regularizing it away: FSQ~\cite{mentzer2024fsq} quantizes each embedding dimension independently, and SimVQ~\cite{zhu2024simvq} reparameterizes the codebook through a single linear layer over a shared basis; we study balance regularization specifically because it is what GENIUS's released architecture uses, not because it is the only remedy to collapse. Concurrent recommendation work reports a related ``hourglass'' code concentration~\cite{kuai2024hourglass} and decouples reconstruction from distribution-matching \emph{within one item stream}~\cite{wang2026decoupledrq}; we instead decouple \emph{across} a cross-modal retriever's two target modalities, failing for a different, architecture-specific reason (Section~\ref{sec:direction-aware}). LETTER~\cite{wang2024letter} adds a diversity loss to RQ-VAE codebooks for generative recommendation -- the closest prior use of our exact intervention -- as one term in a joint objective for a single-modality identifier, not studying its effect on a codebook serving two target modalities. None of this prior work reports whether balance regularization helps or hurts \emph{retrieval}, or whether its effect is stable across regimes -- the gap this paper closes.
